\documentclass[a4paper,12pt]{article}
\usepackage[utf8]{inputenc}
\usepackage[T1]{fontenc}
\usepackage[ngerman]{babel}
\usepackage{amsmath}
\usepackage{amssymb}
\usepackage{enumitem}
\usepackage{geometry}
\geometry{a4paper, margin=2.5cm}

\title{Einsendeaufgaben -- Aufgabe 2.2\\[0.5em]
\large 61211 Analysis}
\author{}
\date{}

\begin{document}
\maketitle

Sei $a \in \mathbb{Q}$ beliebig.

\underline{Fall 1: $a^2 > 2$}

\underline{Fall 1.1: $a < 0$}

Dann gilt $a < -\sqrt{2}$. Es existiert ein $\varepsilon > 0$ mit $a < a + \varepsilon < -\sqrt{2}$. Es existiert also eine $U_\varepsilon(a)$ mit $x^2 > 2$ für alle $x \in U_\varepsilon(a)$.
Nach Definition von $f$ folgt:
\[
    f(x) = 1 \quad \forall x \in U_\varepsilon(a) \cap \mathbb{Q}
\]

Die Restriktion $f|_{U_\varepsilon(a) \cap \mathbb{Q}}$ ist demnach stetig.
Da also $f|_{U_\varepsilon(a) \cap \mathbb{Q}}$ stetig ist, folgt aus dem Satz 2.3.4(i), dass $f$ in $a$ stetig ist.

\underline{Fall 1.2: $a \geq 0$ und Fall 2: $a^2 < 2$}

Nach dem gleichen Prinzip ist $f$ in $a$ stetig.

\bigskip

\begin{enumerate}[label=\alph*)]
    \setcounter{enumi}{3}
    \item
    Wir beweisen die Stetigkeit von $g$ und $h$ mithilfe des Folgenkriteriums.

    Wir untersuchen $g$ mit einem beliebigen $a \in \mathbb{R}$. Sei $y \in \mathbb{R}$ und $(y_n) \to y$ beliebig.

    Zu zeigen ist $g(y_n) \to g(y)$, was äquivalent zu $f(a, y_n) \to f(a, y)$.

    Wir definieren die Folge $(z_n)$ mit
    \[
        z_n := (a, y_n)
    \]

    Da $z_{n1} \to a$ und $z_{n2} \to y$, folgt $z_n \to (a, y)$. Da $f$ stetig ist, konvergiert $(f(z_n))$ gegen $f(a, y)$, was äquivalent zu $g(y_n) \to g(y)$ ist.

    Also ist $g$ in $y$ stetig. Da $y$ beliebig war, ist $g$ auf ganz $\mathbb{R}$ stetig.

    Die Argumentation für $h$ ist ähnlich.

\end{enumerate}

\end{document}
